\chapter*{머리말}
\addcontentsline{toc}{chapter}{머리말}

Volume 2에서는 대수적 체확장, 분해체, 분리성과 정규성, 유한체, 그리고 유한 갈루아 대응을 준비했다. Volume 3에서는 이 사전을 실제 방정식에 적용한다. 출발점은 단순하다.

\begin{quote}
다항식의 갈루아군은 근들을 아무렇게나 섞는 군이 아니라, 근들 사이의 모든 대수적 관계를 보존하는 순열들의 군이다.
\end{quote}

첫 장에서는 분리다항식의 분해체에서 갈루아군이 대칭군의 부분군으로 들어가는 이유를 다시 정리하고, 기약성과 추이성의 관계를 계산 도구로 바꾼다. 이차방정식, 삼차방정식, 특별한 사차방정식과 일반 방정식을 차례로 살펴본다. 삼차다항식에서는 판별식의 제곱근이 갈루아군을 $A_3$와 $S_3$ 사이에서 구별하는 첫 불변량으로 등장한다.

이 권의 원칙은 앞 권과 같다.

\begin{itemize}
  \item 계산 이전에 어떤 군론적 정보가 필요한지 밝힌다.
  \item 근의 공식보다 분해체와 자동동형을 우선한다.
  \item 예제에서는 기약성, 분리성, 분해체 차수와 군구조를 독립적으로 확인한다.
  \item 정리의 가정이 빠지면 어떤 결론이 무너지는지 반례로 확인한다.
  \item 문제는 계산, 증명, 구조의 세 층위로 나누고 핵심 문제에는 상세 풀이를 제공한다.
\end{itemize}

Volume 3 Draft 0.1은 방정식의 갈루아군을 계산하는 첫 원리를 다룬다. 뒤의 장에서는 판별식과 종결식, 원분체와 순환확장, 노름과 대각합, 근호확장, 가해군과 근호해법으로 진행한다.
